\cotefile{19}
\scanpage{50}


% letterhead
\begin{center}
INSTITUT\\
DES HAUTES ÉTUDES SCIENTIFIQUES\\
5, ROND-POINT BUGEAUD, PARIS XVI\textsuperscript{e}\\
KLÉBER 48-56
\end{center}

% paper watermark; the repeated legend is cropped at the edges
\unsure{RJON|RION} EXTRA STRONG \unsure{BAR|BA}

% archive sheet number: 49
\scanpage{51}


\uline{Lemme} Soit $x\in\mathrm{Ob}\,\underline C$, et soit $\bar x$ une enveloppe injective dans $\underline A$, alors $i^!(\bar x)$ est une enveloppe injective de $x$ dans $\underline C$, [donc égale à $x$ si $x$ est Inj. dans $\underline C$]

\uline{Cor.} Les objets injectifs\inserted{ de $\underline C$}, \deleted{\unsure{sont ceux}} \deleted{de la forme} \inserted{ceux isom. à des objets de la forme} $i^!(y)$, $y$ un objet injectif de $\underline A$.

\uline{Cor}\quad $(\mathrm i)\Longleftrightarrow(\mathrm{ii})\Longleftrightarrow(\mathrm{iii})$

\deleted{\unsure{or $i$|donc $i$}}

$(\unsure{\mathrm{iii}|\mathrm{ii}})\Rightarrow(\mathrm{iv})$ \unsure{car} si $x$ est injectif, \unsure{un|son} \uline{supplémentaire} de $i_*i^!(x)$ dans $x$ est \deleted{\unsure{de}} \deleted{injectif} \inserted{$\underline C$-fermé} par loc. cit cor. 2,

% archive sheet number: 50
% The faint vertical paper watermark shows through from the sheet of scan 50.
\scanpage{52}


% letterhead
\begin{center}
INSTITUT\\
DES HAUTES ÉTUDES SCIENTIFIQUES\\
5, ROND-POINT BUGEAUD, PARIS XVI\textsuperscript{e}\\
KLÉBER 48-56
\end{center}

% paper watermark
\unsure{RJON|RION} EXTRA STRONG \unsure{BAR|BA}

% archive sheet number: 51
\scanpage{53}


\margin{$\operatorname{Hom}(x,Sy)=\operatorname{Hom}(Tx,y)$}

% letterhead
\begin{center}
INSTITUT\\
DES HAUTES ÉTUDES SCIENTIFIQUES\\
5, ROND-POINT BUGEAUD, PARIS XVI\textsuperscript{e}\\
KLÉBER 48-56
\end{center}

\uline{Prop}

\diagram{\begin{tikzcd}[ampersand replacement=\&,column sep=large]
A\arrow[r,shift left,"T"] \& B\arrow[l,shift left,dashed,"S"]
\end{tikzcd}}

$T$ foncteur ``passage au quotient'' \unsure{par|pour} une catégorie de fractions.

Soit $x\in\mathrm{Ob}\,A$, $y=T(x)$. Propriétés équivalentes

a) $x$ est libre au dessus de $y$, i.e. $\forall\,x'\in\mathrm{Ob}\,A$,
\[
\operatorname{Hom}_{A}(x',x)\xrightarrow{\sim}\operatorname{Hom}_{B}(Tx',Tx).
\]

b) Si $x'\unsure{\to|-}x''$ dans $A$ est un ``$\deleted{T}C$ isomorphisme'', (i.e. $Tx'\unsure{-|\to}Tx''$ un isomorphisme) alors
\[
\operatorname{Hom}_{A}(x'',x)\unsure{\to|-}\operatorname{Hom}_{A}(x',x)
\]
est une bijection.

c) [Si $B=A/C$, $A$ catégorie abélienne, $C$ sous-catégorie abélienne épaisse] Tout sous-objet de $x$ appartenant à $\mathrm{Ob}\,C$ est \uline{nul}, et $x$ est facteur direct dans tt $y\supset x$ tel que $y/x\in\mathrm{Ob}\,C$

d) [Si l'adjoint $S$ de $T$ existe]
\[
x\xrightarrow{\unsure{\mathrm{can}|\infty|\sim}}STx\quad\text{est un isomorphisme}
\]

e) \inserted{[Si l'adjoint $S$ de $T$ existe]} $x$ est dans l'image essentielle de $S$, i.e. isomorphe à un objet de la forme $Tz\ [z\in\mathrm{Ob}\,C]$

% archive sheet number: 52
\scanpage{54}


\deleted{\uline{Corollaire} Soit \deleted{$x\in\mathrm{Ob}\,A$, \unsure{$y$ $y\in$} et} $x\to x'$ une flèche de $A$. Conditions équivalentes a) $Tx\xrightarrow{\sim}Tx'$, \uline{et} $x'$ est libre sur $Tx'$. \unsure{(cf prop. préc.)|Cf. prop. préc.}}

\uline{Corollaire} 1 Pour que le foncteur adjoint $S$ existe, il f et s. que pour tt $x\in\mathrm{Ob}\,A$, existe une flèche $x\to x'$ telle que (i) $Tx\unsure{-|\to}Tx'$ soit un isom. (ii) $x'$ est libre sur $Tx'$, i.e. satisfait aux conditions de la prop.

D'ailleurs, Gabriel prouve

\uline{Corollaire 2} \deleted{Si $\underline C$} Dans le cas $\underline A$ abélien, $\underline B=\underline A/\underline C$ ($\unsure{\underline C|C}$ sous-catégorie épaisse), pour que $\deleted{\unsure{S(x)|ST}}\inserted{S}$ existe, il f et s. que

a) $\forall\unsure{x\in\mathrm{Ob}\,C}$, existe un plus-grand sous-objet \inserted{$x''$} de $x$ qui soit $\in\mathrm{Ob}\,C$

% slanted
\margin{Il n'est \unsure{pas} possible de donner ceci comme critère pour l'existence de $STx$, $x$ \uline{fixé}.}

b) \inserted{si $x''=0$,} Il existe un monomorphisme de $x\deleted{\unsure{/x''}}$ dans un objet ``$C$-fermé''.

Si $\underline C$ est à $\varinjlim$ exactes et avec générateurs, la condition a) suffit, et signifie aussi que $\underline C$ est \uline{\uline{stable} par $\varinjlim$}. \uline{Alors} $T$ transforme injectifs en injectifs.

% letterhead
\begin{center}
INSTITUT\\
DES HAUTES ÉTUDES SCIENTIFIQUES\\
5, ROND-POINT BUGEAUD, PARIS XVI\textsuperscript{e}\\
KLÉBER 48-56
\end{center}

% archive sheet number: 53
\scanpage{55}


\unsure{Théorèmes|Théorème} de descente de M. Artin - Cartier

11 p

% archive sheet number: 54
\scanpage{57}


\mspage{1}

\uline{Th. de M. Artin et \unsure{P|R}. Cartier}

Soient $C,C',C''$ des catégories
\diagram{\begin{tikzcd}[ampersand replacement=\&,column sep=large,execute at end picture={\draw[thin] (prod.west) to[out=90,in=180] (prod.north) to[out=0,in=90] (prod.east) to[out=-90,in=0] (prod.south) to[out=180,in=-90] (prod.west) -- cycle;}]
f^*=(\deleted{\unsure{j}}i^*,j^*)\colon C \arrow[r] \& |[alias=prod]| C'\times C'' \& {}=D.\quad\text{un foncteur,}
\arrow[from=1-2,to=1-1,start anchor=south east,end anchor=south east,no head,to path={-- ++(0,-.14) -| (\tikztotarget)}]
\end{tikzcd}}

On suppose que $i^*,j^*$ ont des adjoints $i_*,j_*:$
\[
\operatorname{Hom}(F,i_*G)=\operatorname{Hom}(i^*F,G),\quad
\operatorname{Hom}(F,j_*H)=\operatorname{Hom}(j^*F,H)
\]
donc que $f_{\deleted{\unsure{*}}}^{*}$ a un adjoint $f_*$
\[
f_*(G,H)\simeq i_*(G)\times j_*(H)
\]
\[
\begin{aligned}
[\,\operatorname{Hom}(F,f_*(G,H))&\simeq\operatorname{Hom}(F,i_*(G))\times\operatorname{Hom}(F,j_*H)\\
&\simeq\operatorname{Hom}(i^*F,G)\times\operatorname{Hom}(j^*F,H)\\
&\simeq\operatorname{Hom}(f^*F,(G,H))\,]
\end{aligned}
\]

\margin{$+$}

On suppose de plus

\[
\begin{aligned}
\text{\textcircled{a}}\quad \tikz[baseline=(p.base)]{\node[inner xsep=7pt,inner ysep=3pt] (p) {$\deleted{\unsure{j}}i^*\deleted{\unsure{j}}i_*\deleted{\unsure{H}}G\xleftarrow{\sim}G$};\draw[thin] (p.west) to[out=90,in=180] (p.north) to[out=0,in=90] (p.east) to[out=-90,in=0] (p.south) to[out=180,in=-90] (p.west) -- cycle;}&\qquad\text{i.e. $i_*$ est pleinement fidèle}\\
\tikz[baseline=(p.base)]{\node[inner xsep=7pt,inner ysep=3pt] (p) {$j^*j_*H\xleftarrow{\sim}H$};\draw[thin] (p.west) to[out=90,in=180] (p.north) to[out=0,in=90] (p.east) to[out=-90,in=0] (p.south) to[out=180,in=-90] (p.west) -- cycle;}&\qquad\text{i.e. $j_*$ est pleinement fidèle}
\end{aligned}
\]
% The two formulas are separately ovalled in pencil.

\margin{$\times$}

\textcircled{b} $i^*j_*(H)$ est objet final de $C'$ pour tt $H\in\mathrm{Ob}\,C''$

[Alors $f^*f_*(G,H)\simeq(G,\varphi(G)\times H)$, et \deleted{\unsure{alors}} ``donnée de Cartier'' sur $(G,H)$ rel. à $f^*$ est \unsure{donc|d'un|dans} morphisme $(G,\varphi(G)\times H)\leftarrow(G,H)$ de la forme $(\mathrm{id}_G,(u,\mathrm{id}_H))$, où $u:H\to\varphi(G)$ est arbitraire].

% slanted
\margin{N.B. on pose $\varphi=j^*i_*$ [\unsure{commute} aux noyaux de couples en vertu de c]}

\textcircled{c} \inserted{les noyaux de couples existent dans $C$, et} $f^*$ commute aux noyaux de couples [i.e. $i^*,j^*$ y commutent]

% archive sheet number: 56
\scanpage{58}


\margin{\figure{\begin{tikzpicture}\draw (0,0) circle (.045) -- (.35,-.38);\end{tikzpicture}}}\textcircled{d} $f^*$ est conservatif, i.e. si $\lambda\in\mathrm{Fl}\,C$ est tel que $i^*\lambda$ et $j^*\lambda$ soient des isom., alors $\lambda$ est un isom.

Sous ces conditions, le foncteur naturel
\diagram{\begin{tikzcd}[ampersand replacement=\&] C\arrow[r,"\theta"] \& K \end{tikzcd}}
[où $K$ est la \deleted{\unsure{cat. des}} catégorie des triples $(G,H,u)$, $G\in\mathrm{Ob}\,C'$, $H\in\mathrm{Ob}\,C''$, $u\in\operatorname{Hom}(H,\varphi(G))$] est une \uline{équivalence de catégories}.

\uline{Ex} \deleted{\unsure{Espa}}
\[
\begin{aligned}
C&=\text{faisceaux sur }X\\
C'&\unsure{\,\cdot\,|=}\text{faisceaux sur } U\subset X\quad(U\text{ ouvert})\\
C''&=\text{faisceaux sur }Y\subset X\quad(Y=X-U)
\end{aligned}
\]

\margin{\figure{\begin{tikzpicture}\draw (0,0) circle (.035) -- (-.17,-.47);\end{tikzpicture}}}équivaut (en présence des autres axiomes) \deleted{\unsure{au fait que}} \inserted{à la condition} \deleted{diagramme} d') pour tout $F\in\mathrm{Ob}\,C$, le diagramme
\diagram{\begin{tikzcd}[ampersand replacement=\&]
F\arrow[r]\arrow[d] \& i_*i^*F\arrow[d]\\
j_*j^*F\arrow[r] \& j_*j^*i_*i^*F
\end{tikzcd}}
\diagram{\begin{tikzcd}[ampersand replacement=\&]
\text{( i.e.} \& F\arrow[r]\arrow[d] \& i_*(G)\arrow[d] \& {}\\
{} \& j_*(H)\arrow[r,no head,"\unsure{-|\to}"{marking}] \& j_*(\varphi(G)) \& \text{)}
\end{tikzcd}}
est cartésien

% archive sheet number: 57
\scanpage{59}


\mspage{2}

La situation
\diagram{\begin{tikzcd}[ampersand replacement=\&,column sep=large,row sep=large]
\biggl\{\, C''\arrow[r,"j_*"] \& C\arrow[l,shift right=5,"j^*"']\arrow[l,shift left=5,"j^!"]\arrow[r,"\deleted{\unsure{i_*}}i^*"] \& C'\arrow[l,shift right=5,"i_!"']\arrow[l,shift left=5,"i_*"]\\
H \& F \& G
\end{tikzcd}}

(\uline{catégories abéliennes}, \uline{foncteurs additifs})

\[
\left\{\begin{aligned}
\operatorname{Hom}(F,j_*H)&=\operatorname{Hom}(j^*F,H)\\
\operatorname{Hom}(j_*H,F)&\simeq\operatorname{Hom}(H,j^!F)\\
\operatorname{Hom}(F,i_*G)&\simeq\operatorname{Hom}(i^*F,G)\\
\operatorname{Hom}(i_!\deleted{\unsure{F}}G,F)&\simeq\operatorname{Hom}(G,i^*F)
\end{aligned}\right.
\]

\[
\begin{aligned}
j_!=j_*,\,j^*,\,i^*=i^!,\,i_!&\quad\text{exacts}\\
j^!,\,i_*&\quad\text{exacts à gauche}
\end{aligned}
\]

\[
\varphi=j^*i_*:\ C''\longleftarrow C'\qquad\varphi\ \text{exact à gauche}
\]

\[
\left\{\begin{aligned}
j_*j^!F&=\underline{\Gamma}_Y F& j_*j^*F&=F^Y\\
i_*i^*F&=\underline{\Gamma}_U F& i_!i^*F&=F^U
\end{aligned}\right.
\]

\[
\left\{\begin{aligned}
i^*j_*&=0&\quad &\times\\
j^!i_*&=0&&\times\\
j^*i_!&=0&&\times\\
j^!i_!&=0
\end{aligned}\right.
\]

\[
j^*\longrightarrow\varphi i^*
\]

\uline{Prop. 1} Partons de $C''\xrightarrow{j_*}C\xrightarrow{i^*}C'$, où $j_*$ est pleinement fidèle, $i^*$ est un passage au quotient exact, et $j_*\simeq\operatorname{Ker}i^*$. \inserted{donc $i^*j_*=0$,} Alors

(i) Si $i_!$ existe\inserted{ [néc. pl. fid]}, alors $j^*$ existe\inserted{ [néc. un passage au quotient]}, et on a la suite exacte

\diagram{\begin{tikzcd}[ampersand replacement=\&,column sep=small]
C''\arrow[r,shift right,"j_*"'] \& C'\arrow[l,shift right,"j^*"']\arrow[r,shift right,"i^*"'] \& C''\arrow[l,shift right,"i_!"']
\end{tikzcd}}

\[
\boxed{i_!i^*F\longrightarrow F\longrightarrow j_*j^*F\longrightarrow 0}
\]

\[
\boxed{\deleted{\unsure{j^*i_!=0}}}
\]

De plus
\[
\left\{\begin{aligned}
j^*j_*&\overset{\sim}{\longleftrightarrow}\mathrm{id}_{C''}\\
i^*i_!&\xrightarrow{\sim}\mathrm{id}_{C'}\\
j^*i_!&\simeq 0_{C''}
\end{aligned}\right.
\]

(i bis) Si $i_*$ existe, alors $j^!$ existe, et on a
\[
i_*i^*F\longleftarrow F\longleftarrow j_*j^!F\longleftarrow 0
\]

De plus
\[
\left\{\begin{aligned}
j^!j_*&\xleftarrow{\sim}\mathrm{id}_{C''}\\
i^*i_*&\xrightarrow{\sim}\mathrm{id}_{C'}\\
j^!i_*&\simeq 0_{C'}
\end{aligned}\right.
\]

\margin{\uline{Question} Si $j^*$ existe est-il vrai que $i_!$ existe ?}

(N.B. En général, on ne peut rajouter des zéros\ldots)

% archive sheet number: 58
\scanpage{60}


\mspage{3}

L'énoncé (i bis) est dual de (i).

\uline{Corollaire 1} \inserted{Sous les cond. de prop. 1.} Supposons que $i_!$ existe (condition (i) de la proposition 1). Les conditions suivantes sont équivalentes

a) $j^*$ exact, $i_!\simeq\operatorname{Ker}j^*$ [à fortiori, $i_!$ \uline{est exact}\deleted{\unsure{et pl. fid}}]

b) $i_!$ a une image \inserted{essentielle} épaisse [épaisse à g. suffit]

\deleted{\unsure{b}}c) $i_!i^*(F)\unsure{\to|-}F$ est injectif pour tt $F$ i.e. la suite
\[
0\longrightarrow\underset{\substack{\shortparallel\\F^U}}{i_!i^*F}
\longrightarrow F\longrightarrow\underset{\substack{\shortparallel\\F^Y}}{j_*j^*F}
\longrightarrow 0\qquad\text{est exacte}
\]

d) $j^*$ \uline{est exact}

\deleted{Alors $i_!\simeq\operatorname{Ker}j^*$, à fortiori $i_!$ est exact}

a) $\Rightarrow$ b) trivial,

b) $\Rightarrow$ c)
\[
\left\{\begin{aligned}
&\text{Car }0\to N\to i_!i^*F\to F \overset{i^*\text{ exact}}{\Rightarrow}0\to i^*N\to\underbrace{i^*i_!}\,\underbrace{i^*F}\to\underbrace{i^*F}\quad\substack{\text{dernière fl. un isom.}\\\text{($i^*$ passage au quotient)}}\\
&\text{donc }i^*(N)=0\ \deleted{\unsure{\text{donc}}},\ \text{or } N=i_!(M)\ \text{par b)},\quad i^*i_!(M)=0\Rightarrow M=0,\Rightarrow N=0
\end{aligned}\right.
\]

c) $\Rightarrow$ a) \deleted{\unsure{$F\to F'\to F''\to 0$}} Supposons \uline{$0\to F\to F'\to F''\to0$} d'où

\deleted{a) $\Rightarrow$ d) trivial}\quad\deleted{\unsure{d) $\Rightarrow$ a)|d) $\Rightarrow$ b)}}

\diagram{\begin{tikzcd}[ampersand replacement=\&,column sep=small,row sep=small]
{} \& {} \& \arrow[d] {} \& {} \& {}\\
\deleted{0}\arrow[r] \& i_!i^*F\arrow[r]\arrow[d] \& F\arrow[r]\arrow[d] \& j_*j^*F\arrow[r]\arrow[d] \& 0\\
\deleted{\unsure{0\to}} \& i_!i^*F'\arrow[r]\arrow[d] \& F'\arrow[r]\arrow[d] \& j_*j^*F'\arrow[r] \& 0\\
0\arrow[r] \& i_!i^*F''\arrow[r] \& F''
\end{tikzcd}}

et utilisant l'injectivité de $i_!i^*F''\to F''$, on trouve $j_*j^*F\to j_*j^*F'$ est injectif, donc $j^*F\to j^*F'$ inj. d'où $j^*$ ex. (car il est déjà exact à droite).

Soit $F=i_!G$, alors $i_!i^*F\to F$ est un isom [$i^*$ passage au quotient] donc $j_*j^*F=0$ \deleted{\unsure{est un isom}} par la suite exacte (i)\inserted{ déjà vu dans prop. 1.} donc $j^*F=0$. Réciproquement, si $j^*F\unsure{=0|-0}$, alors par c), on trouve $i_!i^*F\xrightarrow{\sim}F$, i.e. $F\unsure{\in|\equiv}\operatorname{Im}\,\mathrm{ess}\,i_!$

% archive sheet number: 59
\scanpage{61}


\mspage{4}

a) $\Rightarrow$ d) trivial

d) $\Rightarrow$ c) Car \deleted{\unsure{d) si on}} $0\to N\to i_!i^*F\to F$ on aura \inserted{par d)} $0\to j^*N\to\underset{\substack{\shortparallel\\0}}{j^*i_!}i^*F$ \deleted{\unsure{exacte}} d'où $j^*N=0$, or on a \deleted{\unsure{$i^*N=0$}} aussi $i^*N=0$ [$0\to i^*N\to i^*i_!i^*F\to i^*F$] donc $N=j_*\deleted{\unsure{N|M}}H$, $H\in\mathrm{Ob}\,C''$, et comme $H=j^*\deleted{\unsure{N}}j_*H\simeq j^*N=0$, on aura $N=0$, cqfd

\margin{Cor. 2 \rule{1.2cm}{0.3pt}}

\uline{Corollaire 1} bis dual. (Les conditions \unsure{du dit|duales|ci-dessus} ne seront pas vérifiées dans les applications en vue, en particulier $j^!$ ne sera pas exact, $\operatorname{Im}i_*$ ne sera pas épaisse à droite, $F\to i_*i^*F$ ne sera pas en général surjectif].

\uline{\unsure{Remarque|Remarques}} 1 \deleted{\unsure{Si}} Les conditions du cor. 1 ne \uline{résultent pas} du fait que $i_!$ est exact, comme on voit dans la situation duale pour la localisation des modules \ldots

\uline{Remarque} 2 Sous les conditions équiv. du Cor. 1, on a la situation
\diagram{\begin{tikzcd}[ampersand replacement=\&,column sep=large]
C''\arrow[r,shift right,"j_*"] \& C\arrow[l,shift right,"j^*"']\arrow[r,shift right,"i^*"] \& C'\arrow[l,shift right,"i_!"']
\end{tikzcd}}
parfaitement \uline{symétrique} en $C',C''\deleted{\unsure{\, '|\circ}}:$
\[
\left\{\begin{aligned}
&i^*\ \text{passage au }\deleted{\text{quotient}}\ \text{exact de noyau } i_!\\
&j^*\ \text{passage au quotient exact de noyau }i_!
\end{aligned}\right.
\]
$j_*$ et $j^*$ adjoints, $i_{\deleted{\unsure{!}}}^*$ et $i_{\smash{!}}^{\deleted{\unsure{*}}}$ adjoints

De façon précise, la situation
\diagram{\begin{tikzcd}[ampersand replacement=\&,column sep=large]
C^{\prime o}\arrow[r,shift right,"i_!^{o}"] \& C^{o}\arrow[l,shift right,"i^{*o}"']\arrow[r,shift right,"j^{*o}"] \& C^{\prime\prime o}\arrow[l,shift right,"j_*^{o}"']
\end{tikzcd}}
satisfait les \deleted{\unsure{co}} \unsure{mêmes|mes|aux} hyp. que la situation de départ Cette symétrie va être détruite dans ce qui suit.

% archive sheet number: 60
\scanpage{62}


\mspage{5}

\uline{\deleted{\unsure{Prop}} \unsure{Cor|Corol|Cor.} 2} Sous les conditions \unsure{de|du} prop. 1, conditions équivalentes :

a) $j^*$ existe \inserted{[néc. passage au quotient]} et est exact

b) $i_!$ existe \inserted{[néc. pl. fid]}, et a une image \inserted{essentielle} épaisse

b') $i_!$ existe, et $i_!i^*F\unsure{-|\to}F$ est injective pour tout $F$.

Dans ce cas, $i_!\simeq\operatorname{Ker}j^*$, $j^*\simeq\operatorname{Coker}i_!$, à fortiori $i_!$ exact

[N.B. Mais il ne \uline{suffit} pas que $j^*,i_!$ existent, $i_!$ exact \ldots]

La prop. 1 et son corollaire prouvent que a) implique b), b').

% Faint paper watermark, visible through the sheet: XTRA STRONG.
% archive sheet number: 61
\scanpage{63}


\mspage{5 bis}

\diagram{\begin{tikzcd}[ampersand replacement=\&]
\deleted{C}\arrow[r] \& \deleted{D}
\end{tikzcd}}

\diagram{\begin{tikzcd}[ampersand replacement=\&,column sep=small]
\deleted{C''}\arrow[r,no head,shift right,"\deleted{j_*}"'] \& \deleted{C}\arrow[l,shift right,"\deleted{j^*}"']\arrow[r,no head,"\deleted{i^*}"'] \& \deleted{C'}\\
{} \& \deleted{i_!i^*F} \& {}
\end{tikzcd}}

\deleted{\uline{Cor} Equivalent conditions, for $C''\xrightarrow{j_*}C'\xrightarrow{i^*}C''$ as above}

\deleted{a) $i_!$ exists, and has thick image [or : $i_!i^*F\to F$ injective]}

\deleted{b) $j^*$ exists and is \uline{exact}}

\margin{\deleted{[N.B. even if $i_!$ exists and is exact, so that $j^*$ exists, \unsure{we have seen|we don't see} that $j^*$ may not be exact]}}

\deleted{If so, $i_!\simeq\operatorname{Ker}j^*$, $j^*\simeq\operatorname{Coker}i_!$}

We have seen already a) $\Rightarrow$ b), let us prove b) $\Rightarrow$ a). Therefore, let us define the functor
\[
\tau(F)=\operatorname{Ker}(F\to j_*j^*F)
\]
we see that $\tau$ is an \deleted{\unsure{an}} exact functor [thanks to the fact that $j_*j^*F$ is exact] Zero on \unsure{objects} of the type $j_*F'$, \deleted{\unsure{there}} hence we have
\[
\tau(F)\simeq i_!i^*F
\]
with a well defined \inserted{exact} functor $i_!:C''\unsure{-|\to}C$. Now \inserted{to define}
\[
\operatorname{Hom}(i_!\unsure{G|G''},\deleted{\unsure{j^*}}F)\unsure{.}\simeq\operatorname{Hom}(G,i^*F),\qquad\text{we need}
\]
to define only
\[
\operatorname{Hom}(i_!i^*F',F)\simeq\operatorname{Hom}(i^*F',i^*F)
\]
\[
\text{i.e.}\qquad\operatorname{Hom}(\tau F',F)\simeq\operatorname{Hom}(i^*F,i^*F)
\]

But \uline{as} $j^*F$ \uline{is the biggest} \deleted{\unsure{negligeable|negligible}} \uline{quotient} of $F$ \uline{in} \deleted{\unsure{$j_*C'$}}, $\operatorname{Im}j_*=\operatorname{Ker}i^*$, this comes out immediately, by the construction of a quotient category $C/C''$

% slanted
\margin{to give us another formulation for the \unsure{construction|criterion} of $j^*$, as well as}

\uline{Question} \inserted{\unsure{Does|does}} the existence of $j^*$ (exact or not) \deleted{does} \deleted{\unsure{it}} imply the existence of $i_!$ ?

% archive sheet number: 62
\scanpage{64}


\[
\begin{aligned}
j^*i_!&=0 &i_!&=\operatorname{Ker}j^* &j^*&=\operatorname{Coker}i_!\\
i^*j_*&=0 &j_*&=\operatorname{Ker}i^* &i^*&=\operatorname{Coker}j_*
\end{aligned}
\]

\[
\begin{aligned}
&0\unsure{-|\to} N\to i_!i^*F\to\unsure{F|B}\\
&0\to i^*N\to i^*i_!i^*F\to i^*F\\
&0\to i_!i^*N\to i_!i^*i_!i^*F\to i_!i^*F
\end{aligned}
\]

\diagram{\begin{tikzcd}[ampersand replacement=\&,column sep=large,row sep=small]
\deleted{C''}\arrow[r,shift right,"\deleted{j_*=j_!}"'] \& \deleted{C}\arrow[l,shift right,"\deleted{j^*}"']\arrow[r,shift right,"\deleted{i^*=i^!}"'] \& \deleted{C'}\arrow[l,shift right,"\deleted{i_!}"']\\
{} \& {}\arrow[l,"\deleted{j^!}"] \& {}\arrow[l,"\deleted{i_*}"]
\end{tikzcd}}

\deleted{Equivalent conditions on pair $i^*,i_!$}

\deleted{a) $i_!$ is \uline{fully faithful}, \uline{thick image, and $i_!i^*F\unsure{-|\to}F$ injective}}

\deleted{b) $i^*$ \uline{passage to quotient} \uline{that is exact}, and one of the following conditions hold (All \unsure{expressed} by the previous assumption \unsure{b)}; $j_*$ is defined thanks to said assumption, and $j^*$ then must exist as we declared,] and is exact passage to quotient]}

\deleted{$\alpha$) $i_!$ \deleted{\unsure{is}} has thick image}

\deleted{$\beta$) $i_!i^*F\to F$ \unsure{injection|injective} for every $F$}

\deleted{$\unsure{\gamma|\delta}$) $i_!=\operatorname{Ker}j^*$}

\deleted{[ N.B. \unsure{Mixture|Many} of \inserted{\unsure{equiv.}} conditions which are formally distinct )}

\deleted{This situation is essentially autodual. We have then an exact sequence [cf. \uline{Cartier}]}

\[
\boxed{\deleted{0\longrightarrow i_!i^*F\longrightarrow F\longrightarrow j_*j^*F\longrightarrow0}}
\]

\deleted{Under these conditions, is it true that the existence of $i_*$ is essentially equivalent to the existence of $\unsure{j^!|j_!|j?}$ Of course not, it only \uline{implies} the one of $j^!$ by}

\[
\boxed{\deleted{j_*j_!F=\operatorname{Ker}(F\to i_*i^*F)}}
\]

\deleted{$\unsure{i_*G}$}

% archive sheet number: 63
